% Part: first-order-logic
% Chapter: axiomatic-deduction
% Section: provability-propositional

\documentclass[../../../include/open-logic-section]{subfiles}

\begin{document}

\iftag{FOL}
      {\olfileid[fr]{fol}{axd}{ppr}}
      {\olfileid[fr]{pl}{axd}{ppr}}

\olsection{\usetoken{S}{derivability} et connecteurs propositionnels}

\begin{explain}
Nous établissons que la relation de !!{derivability}~$\Proves$ de la
déduction axiomatique est assez puissante pour démontrer certains faits
élémentaires sur les connecteurs propositionnels, par exemple
$!A \land !B \Proves !A$ et
$!A, !A \lif !B \Proves !B$ (modus ponens). Ces faits seront
nécessaires à la preuve du théorème de complétude.
\end{explain}

\begin{prop}\ollabel{prop:provability-land}
\begin{enumerate}
\item \ollabel{prop:provability-land-left}
  $!A \land !B \Proves !A$ et $!A \land !B \Proves !B$;
\item \ollabel{prop:provability-land-right}
  $!A, !B \Proves !A \land !B$.
\end{enumerate}
\end{prop}

\begin{proof}
\begin{enumerate}
\item Ces deux résultats découlent respectivement de
  l'\olref[prp]{ax:land1} et de l'\olref[prp]{ax:land2} par modus ponens.
\item Le résultat découle de l'\olref[prp]{ax:land3} par deux
  applications du modus ponens.
\end{enumerate}
\end{proof}

\begin{prop}\ollabel{prop:provability-lor}
\begin{enumerate}
\item $!A \lor !B, \lnot !A, \lnot !B$ est incohérent.
\item $!A \Proves !A \lor !B$ et $!B \Proves !A \lor !B$.
\end{enumerate}
\end{prop}

\begin{proof}
\begin{enumerate}
\item D'après l'\olref[prp]{ax:lnot2}, nous avons
  $\Proves \lnot !A \lif (!A \lif \lfalse)$ et
  $\Proves \lnot !B \lif (!B \lif \lfalse)$. Le théorème de la
  déduction donne donc
  $\{\lnot !A\} \Proves !A \lif \lfalse$ et
  $\{\lnot !B\} \Proves !B \lif \lfalse$. À partir de
  l'\olref[prp]{ax:lor3}, nous obtenons
  $\{\lnot !A, \lnot !B\} \Proves (!A \lor !B) \lif \lfalse$.
  Le théorème de la déduction donne alors
  $\{!A \lor !B, \lnot !A, \lnot !B\} \Proves \lfalse$.
\item Le premier résultat découle de l'\olref[prp]{ax:lor1}; le second,
  de l'\olref[prp]{ax:lor2}. Dans les deux cas, on applique le modus
  ponens.
\end{enumerate}
\end{proof}

\begin{prop}\ollabel{prop:provability-lif}
\begin{enumerate}
\item \ollabel{prop:provability-lif-left}
  $!A, !A \lif !B \Proves !B$.
\item \ollabel{prop:provability-lif-right}
  $\lnot !A \Proves !A \lif !B$ et $!B \Proves !A \lif !B$.
\end{enumerate}
\end{prop}

\begin{proof}
\begin{enumerate}
\item Nous pouvons donner la !!{derivation} suivante :
  \begin{derivation}
    1. & $!A$ & \Hyp\\
    2. & $!A \lif !B$ & \Hyp\\
    3. & $!B$ & 1, 2, \MP
  \end{derivation}

\item Le premier résultat découle de l'\olref[prp]{ax:lnot2}; le second,
  de l'\olref[prp]{ax:lif1}. Dans les deux cas, on emploie le théorème
  de la déduction.
\end{enumerate}
\end{proof}

\begin{editorial}
Dans la preuve de la première proposition, la source renvoie deux fois à
l'\olref[prp]{ax:land1}; le second renvoi doit viser
l'\olref[prp]{ax:land2}. Dans celle de la disjonction, elle cite
l'\olref[prp]{ax:lnot1}, alors que les deux formules affichées sont des
instances de l'\olref[prp]{ax:lnot2}. Nous avons rectifié ces renvois sans
modifier les formules ni les conclusions.
\end{editorial}

\end{document}
